%%
%%  Annexes
%%
%%  Note: Ne pas modifier la ligne ci-dessous. / Do not modify the following line.
\ifthenelse{\equal{\Langue}{english}}{
	\addcontentsline{toc}{compteur}{APPENDICES}
}{
	\addcontentsline{toc}{compteur}{ANNEXES}
}
%%
\Annexe{Paramètres de simulation communs aux trois articles}\label{ann:params}

Le Tableau~\ref{tab:common-params} résume les paramètres de simulation communs utilisés dans les trois articles de cette thèse.

\begin{table}[htb]
\centering
\caption{Paramètres de simulation communs}
\label{tab:common-params}
\begin{tabular}{lll}
\toprule
\textbf{Paramètre} & \textbf{Valeur} & \textbf{Articles} \\
\midrule
Nombre de n\oe{}uds ($N$) & 20--100 (Art.~1), 200 (Art.~2) & 1, 2 \\
Dimension du terrain & $200 \times 200$\,m$^2$ & 1, 2 \\
Position de la BS & $(100, 250)$ & 1, 2 \\
Énergie initiale ($E_0$) & 0.5\,J & 1, 2 \\
$E_{\text{elec}}$ & 50\,nJ/bit & 1, 2 \\
$\varepsilon_{\text{fs}}$ & 10\,pJ/bit/m$^2$ & 1, 2 \\
$\varepsilon_{\text{mp}}$ & 0.0013\,pJ/bit/m$^4$ & 1, 2 \\
$d_0$ (distance seuil) & 87.7\,m & 1, 2 \\
Taille des paquets & 4\,000\,bits & 1, 2 \\
Nombre de tours & 5\,000 & 1, 2 \\
Tour de catastrophe & 2\,000 & 2 \\
Graines aléatoires & 3 (Art.~1), 30 (Art.~2), 20 (Art.~3) & 1, 2, 3 \\
\bottomrule
\end{tabular}
\end{table}


\Annexe{Architectures des réseaux de neurones}\label{ann:architectures}

\section*{Article 1~: GAPF}

L'architecture GAPF comprend 1\,473 paramètres répartis comme suit~:

\begin{table}[htb]
\centering
\caption{Décomposition des paramètres de GAPF}
\label{tab:gapf-params}
\begin{tabular}{lrr}
\toprule
\textbf{Composant} & \textbf{Dimensions} & \textbf{Paramètres} \\
\midrule
GCN couche 1 ($\mathbf{W}_1$) & $4 \times 16$ & 64 \\
GCN couche 2 ($\mathbf{W}_2$) & $16 \times 16$ & 256 \\
CAS $\mathbf{W}_c$ & $16 \times 64$ & 1\,024 \\
CAS $\mathbf{b}_c$ & $64$ & 64 \\
CAS $\mathbf{w}_2$ & $64 \times 1$ & 64 \\
CAS $b_2$ & $1$ & 1 \\
\midrule
\textbf{Total} & & \textbf{1\,473} \\
\bottomrule
\end{tabular}
\end{table}

\section*{Article 2~: Agent DQN de CALASH}

L'agent Double DQN de CARE utilise l'architecture MLP suivante~:

\begin{table}[htb]
\centering
\caption{Architecture de l'agent DQN de CALASH}
\label{tab:calash-dqn}
\begin{tabular}{lrr}
\toprule
\textbf{Couche} & \textbf{Dimensions} & \textbf{Paramètres} \\
\midrule
Entrée $\to$ Cachée 1 & $9 \times 64$ + 64 & 640 \\
Cachée 1 $\to$ Cachée 2 & $64 \times 32$ + 32 & 2\,080 \\
Cachée 2 $\to$ Sortie & $32 \times 1$ + 1 & 33 \\
\midrule
\textbf{Total} & & \textbf{2\,753} \\
\bottomrule
\end{tabular}
\end{table}

Le vecteur d'état à 9 dimensions encode~: (1)~énergie résiduelle normalisée, (2)~distance au CH, (3)~distance CH$\to$BS, (4)~nombre de membres du cluster, (5)~intensité carbone normalisée, (6)~pression de la file carbone, (7)~ratio de n\oe{}uds vivants, (8)~facteur de cycle de vie, (9)~indicateur de catastrophe.

\section*{Article 3~: CASTER-ZT}

L'architecture CASTER-ZT comprend $\sim$20\,000 paramètres~:

\begin{table}[htb]
\centering
\caption{Décomposition des paramètres de CASTER-ZT}
\label{tab:casterzt-params}
\begin{tabular}{lrr}
\toprule
\textbf{Composant} & \textbf{Architecture} & \textbf{Paramètres} \\
\midrule
Encodeur GraphSAGE & 2 couches, $d{=}64$ & $\sim$8\,400 \\
Politique de récupération & MLP $128 \to 64 \to |\mathcal{A}|$ & $\sim$8\,500 \\
Évaluateur confiance-risque & MLP double-tête & $\sim$3\,000 \\
Bouclier déterministe & 14 seuils scalaires & 14 \\
\midrule
\textbf{Total} & & $\sim$\textbf{20\,000} \\
\bottomrule
\end{tabular}
\end{table}


\Annexe{Protocoles de référence}\label{ann:baselines}

Le Tableau~\ref{tab:baselines} récapitule les protocoles de référence utilisés dans les évaluations expérimentales des trois articles.

\begin{table}[htb]
\centering
\caption{Protocoles de référence par article}
\label{tab:baselines}
\begin{tabular}{llp{3in}}
\toprule
\textbf{Article} & \textbf{Protocole} & \textbf{Description} \\
\midrule
1 & LEACH & Rotation aléatoire des CH \\
1 & EE-LEACH & Élection de CH pondérée par l'énergie \\
1 & ABC-ACO & Hybride métaheuristique \\
1 & Q-Routing & Q-learning tabulaire pour le routage \\
\midrule
2 & LEACH & Rotation aléatoire des CH \\
2 & LEACH-1hop & LEACH avec relais mono-saut \\
2 & EE-LEACH & Élection pondérée par l'énergie \\
2 & ABC-ACO & Hybride métaheuristique \\
2 & EERP & Protocole de routage écoénergétique \\
2 & Q-Routing & Q-learning pour le routage \\
2 & RIS-DRL & RL profond avec surfaces RIS \\
\midrule
3 & CPO-Soft & Optimisation de politique sous contrainte \\
3 & Shield-Binary & Bouclier binaire autoriser/bloquer \\
3 & Agentic-Auto & Politique autonome sans bouclier \\
3 & IF-Trust & Isolation Forest avec confiance \\
\bottomrule
\end{tabular}
\end{table}


